% IEEEtran two-column tweaks for the single-source paper.md build.
\usepackage{tikz}
\usetikzlibrary{arrows.meta,positioning,fit,backgrounds,calc}
\usepackage{booktabs}
\usepackage{array}
\usepackage{amsmath}
\usepackage{amssymb}
\usepackage{fontspec}
% IEEE's look is Times: TeX Gyre Termes as the main font.  DejaVu Mono for
% verbatim so Fig.1's box-drawing glyphs render; shrink code to fit the
% 3.5in column.
\setmainfont{TeX Gyre Termes}
\setmonofont{DejaVu Sans Mono}[Scale=0.80]
\usepackage{fancyvrb}
\usepackage{fvextra}
\fvset{fontsize=\scriptsize,breaklines=true,breakanywhere=true}
\let\oldverbatim\verbatim
\renewcommand{\verbatim}{\scriptsize\oldverbatim}
% Map the Unicode symbols used in prose to LaTeX so the Times font renders
% them (xelatex would otherwise drop missing glyphs silently).
\usepackage{newunicodechar}
\newunicodechar{ψ}{$\psi$}\newunicodechar{ω}{$\omega$}
\newunicodechar{·}{$\cdot$}\newunicodechar{−}{$-$}\newunicodechar{×}{$\times$}
\newunicodechar{→}{$\rightarrow$}\newunicodechar{§}{\S}\newunicodechar{≡}{$\equiv$}
\newunicodechar{≈}{$\approx$}\newunicodechar{≥}{$\geq$}\newunicodechar{≫}{$\gg$}
\newunicodechar{∈}{$\in$}\newunicodechar{⊙}{$\odot$}\newunicodechar{ℓ}{$\ell$}
\newunicodechar{µ}{$\mu$}\newunicodechar{½}{\textonehalf}\newunicodechar{†}{\dag}
\newunicodechar{…}{\ldots}\newunicodechar{¹}{\textsuperscript{1}}
\newunicodechar{²}{\textsuperscript{2}}\newunicodechar{⁰}{\textsuperscript{0}}
\newunicodechar{⁷}{\textsuperscript{7}}\newunicodechar{⁸}{\textsuperscript{8}}
\newunicodechar{⁻}{\textsuperscript{-}}\newunicodechar{₀}{\textsubscript{0}}
\newunicodechar{₁}{\textsubscript{1}}\newunicodechar{₂}{\textsubscript{2}}
\providecommand{\tightlist}{\setlength{\itemsep}{0pt}\setlength{\parskip}{0pt}}
% prose references say "Table 1"; keep captions arabic (IEEEtran defaults to
% Roman numerals)
\renewcommand\thetable{\arabic{table}}
% keep floats near their references: allow double-column floats on the SAME
% page as their source (stock LaTeX defers them to the next page), and relax
% the float limits so a results page may carry several tables
\usepackage{dblfloatfix}
\renewcommand{\topfraction}{0.9}
\renewcommand{\dbltopfraction}{0.9}
\renewcommand{\textfraction}{0.05}
\renewcommand{\floatpagefraction}{0.8}
\setcounter{topnumber}{3}
\setcounter{dbltopnumber}{4}
\setcounter{totalnumber}{5}
